\note{2}{Mon 11 Dec 2023 22:21}{Transforming Hamiltonian to Action Angle Representation}

\section{Hamiltonian Mechanics}

\subsection{First integral}

First integral means independent of $H$ on some coordinate will imply existence of some conserved quantity. For example, if $H(p_1, q_1, p_2, q_2)$ is independent of $q_2$, then
\[
	\dot{p}_2 = \frac{\partial H}{\partial q_2}  = 0
.\] 
Which means the conjugate momentum $p_2$ is conserved. If $H$ is independent of $t$, then $H$ itself is conserved. 

Note that first integral in Hamiltonian dynamics is much simpler than in the Lagrangian, due to the simple form of Hamiltonian equation.




\subsubsection{Canonical transformation}

A canonical transformation 
\[
\begin{aligned}
& p \rightarrow P(q, p) \\
& q \rightarrow Q(q, p) \quad H(q, p, t) \rightarrow K(Q, P, t)
\end{aligned}
\]
from $p, q$ coordinate to $P, Q$ coordinate is given by a transformation function $M(q, Q)$ such that
\[
\sum_{a=1}^N p_a \dot{q}^a-H(q, p, t)=\sum_{a=1}^N P_a \dot{Q}^a-K(Q, P, t)+\frac{d M}{d t}
.\]

\subsubsection{Liouville Theorem}

Recall from vector calculus, that for vector field $f$ and closed region $D$ with smooth boundary $\partial D$, we have
\[
	\oint_{\partial D} f \cdot  \nu dS = \int _D \nabla \cdot f d x
\] 
where $\nu$ is outward normal of boundary.
That is, outward flux can be obtained by integrating the divergence. Let $x = (p, q)$.

Recall also that the total area of $D$ is
\[
	|D| = \int _D d x 
	= \int _D \nabla \cdot (p, 0) d x
	= \oint_{\partial D} (p, 0) \cdot  \nu d S 
	= \oint_{\partial D} p dq
.\] 

Rate of change of phase space volume is
\[
	\frac{dV}{dt} = \oint_{\partial D} \dot{x} \cdot \nu d S
	= \oint_{\partial D} (\nabla \times H) \cdot \nu d S
	= \int_{D} \nabla \cdot  (\nabla \times H) dx
.\] 
since, $\dot{x} = (\frac{\partial H}{\partial q}, - \frac{\partial H}{\partial p} ) = \nabla \times H $. From this we see that, any system described by time-independent Hamiltonian must conserve volume in phase space.

\subsubsection{Conversion to Action-Angle}

The goal is to transform the system described by  $H(p_\alpha, q^\alpha)$ to a system described by $K(J, \theta)$. We may set $H = K$ and let $M(q, \theta) = p q - J \theta$. 

$dM = p d q - J d \theta$ is an exact form. Hence $\oint dM = 0$ along a closed orbit. From this we conclude

\[
	\oint J d \theta = \oint p d q
.\] 
We want $J$ constant along any closed orbit, and any orbit to be described by $2 \pi$-periodicity in $\theta$. Hence, 
\[
	2 \pi J = \oint p d q
.\] 
Now
\[
	\oint_{\partial D} p d q
	=\oint_{\partial D} p \cdot \nu_p d s
	=\iint_D dx = |D|
.\] 

Hence
\[
	J = \frac{1}{2 \pi} \oint_{\partial D} p d q
.\] 
By integrating over $q$, we get $J = J(H)$. Invert to get $H = H(J)$, independent of $\theta$.

Now, angular frequency is given by
\[
	\omega = \dot{\theta} = \frac{\partial H}{\partial J} 
.\] 

Note that $\omega$ is always a constant.

